\documentclass{article}

\usepackage[utf8]{inputenc} % allow utf-8 input
\usepackage[T1]{fontenc}    % use 8-bit T1 fonts
\usepackage{hyperref}       % hyperlinks
\usepackage{url}            % simple URL typesetting
\usepackage{booktabs}       % professional-quality tables
\usepackage{amsfonts}       % blackboard math symbols
\usepackage{nicefrac}       % compact symbols for 1/2, etc.
\usepackage{microtype}      % microtypography


% ---------------------------------------
% 
\usepackage{fullpage}
\usepackage{natbib}
\usepackage{graphicx}

% custom packages
\usepackage{amsthm}
\newtheorem{theorem}{Theorem}
\newtheorem{lemma}{Lemma}
\newtheorem{fact}{Fact}
\newtheorem{assumption}{Assumption}

% custom defintions
\setlength{\parindent}{0pt}

\usepackage{amsmath}
\newcommand\norm[1]{\left\lVert#1\right\rVert}

\usepackage{multirow}
\usepackage[dvipsnames]{xcolor}
\definecolor{first}{HTML}{0077BB}
\newcommand\firstplace[1]{\textcolor{first}{\textbf{#1}}}
\definecolor{second}{HTML}{009988}
\newcommand\secondplace[1]{\textcolor{second}{\underline{\smash{}#1}}}
\definecolor{third}{HTML}{CC3311}
\newcommand\thirdplace[1]{\textcolor{third}{\textit{#1}}}


\title{
    STOR 712 Project Proposal \\
    Can We Improve The Generalizability And Fairness Of Personalized Federated Learning With Learning Rate Scheduling?
}

\author{
    Joseph Lavond and Peter Lin
}



\begin{document}


\maketitle


\section{Background}

\paragraph{Federated Learning (FL)}

Federated Learning (FL) is a potential solution to constructing a shared model (global) from several local data sources that cannot be exchanged or aggregated.
FL is an iterative procedure involving rounds of model improvement until it meets some criteria.
See Figure~\ref{fig: fl} for a graphical representation of the following steps in each iteration of federated training.
These rounds send the current global model to local entities (users) and select some to update the global model. 
The chosen users train their local copy of the model on their data and communicate the resulting model, where the returning models are aggregated to create a new global model. 
The most popular aggregation method of FL is FedAvg~\citep{fedavg}, which sets the new global model to the average of returning models.

\begin{figure}
    \centering
    \includegraphics[width = .7\textwidth]{STOR_712_Project/images/fedavg.png}
    \caption{Single Iteration of Federated Training}
    \label{fig: fl}
\end{figure}


\paragraph{Personalized Federated Learning (pFL).}
When local data are not independent or identically distributed, ensuring good performance across different users with a single shared global model can be challenging.
Recent works \citep{pfl-bench} argue instead that the focus should be on personalized methods that output models adapted to the local data.
The simplest solution involves fine-tuning the shared global model for each user after the federated learning process has terminated.
More complex solutions include learning a model for each user \citep{ditto}, global sharing of a subset of the entire model \citep{fedbn}, and viewing user data as a mixture of distributions \citep{fedem}.
We complain that the above methods require significant increases in computation and memory requirements. 
Can we instead find a global model, without additional computation or memory, that acts as a good initialization for fine-tuning to each user that performs as well as more complex personalized methods?


%
\paragraph{Meta Learning.}
A highly related perspective to personalized federated learning is meta-learning, which attempts to learn a good model initialization that converges quickly for downstream tasks.
Although these tasks are not defined to be closely related, we can think of the personalized federated learning setting as a subset of the field of meta-learning where data heterogeneity is the source of task diversity. 
Several recent attempts have applied meta-learning algorithms to personalized federated learning.
In addition to applying traditional meta-learning algorithms to federated learning \citep{maml}, a newer variant, Reptile~\citep{reptile}, can be shown to be equivalent to FedAvg in the federated learning setting in the case that users have local data of equal size.




%
\section{Proposal}

Consider the following notations.
Let $N$ denote the number of communication rounds of $K$ local update steps for federated learning by users $C = \left[1, \ldots, M \right]$ with local objective functions $F_i$.
Define $\theta_g^{(n)}$ and $\theta_i^{(n, k)}$ as the global model at communication round $n$ and the ith user's copy of that global model after $k$ local update steps with learning rate $\eta$, respectively.
We give the standard update in Equation~\ref{eq: theta}, which gives the FedAvg or FedSGD procedure when $K > 1$ or $K = 1$, respectively.
\begin{align}
    & \theta_i^{(n, k)} = \theta_i^{(n, k - 1)} - \eta \times g_i^{(n, k)} \textrm{ for } k = 1, \ldots, K \label{eq: theta} \textrm{ where } g_i^{(n, k)} = \nabla F_i^{(k)}(\theta_i^{(n, k - 1)}) = \dfrac{1}{\eta} \left( \theta_i^{(n, k - 1)} - \theta_i^{(n, k)} \right) \nonumber
\end{align}

Let $S^{(n)} \subseteq C$ denote the set of users participating in the update of the global model for a given communication round of federated training.
We intend to use the term full user participation for the case $S^{(\cdot)} = C$. 
The global update for full user participation aggregates either the locally updated copy or changes made by local updates. 
See Equation~\ref{eq: global-full-diff}.
\begin{align}
    \theta_g^{(n + 1)} &= \dfrac{1}{M} \sum_{i \in C} \theta_i^{(n, K)} = \theta_g^{(n)} + \dfrac{1}{M} \sum_{i \in C} \left( \theta_i^{(n, K)} - \theta_g^{(n)} \right) \label{eq: global-full-diff}
\end{align}

Depending on the choice of local update, changes made by local training, $\theta_i^{(n, K)} - \theta_g^{(n)}$, can often be expressed as $-\eta \times g_{\textrm{Method}}$.
Equation~\ref{eq: methods} presents the expression for this model difference for FedSGD, FedAvg, and First-Order Model-Agnostic Meta-Learning (FOMAML) described in \cite{maml}.
\begin{align}
    & g_{FedSGD} = g^{(n, 1)}, \qquad g_{FedAvg} = \sum_{j = 1}^K g^{(n, j)}, \qquad g_{FOMAML} = g^{(n, K)}  \label{eq: methods}
\end{align}

In the balanced user data case, FedSGD can be considered an approximation to the situation where all data is aggregated.
Hence, $g^{(\cdot, 1)}$ can be considered highly important for returning a solid initial model, which is necessary for users without sufficient data for fine-tuning.
The other two objectives to simultaneously address are improved personalized models and fast convergence.
Although meta-learning perspectives (FedAvg/Reptile and FOMAML) are designed to adapt well, they place weight on late or all epochs of local training.
We believe that if initial model success and rapid personalization are the main goals for personalized federated learning, the majority of emphasis should be on the early $g_k$. \\

Hence we propose incorporating exponential decay on the gradients of local updates to emphasize early $g_k$. 
We believe intuitively that, with a greater relative emphasis on $g_1$, our proposal may result in better initial accuracy than FedAvg or meta-learning algorithms. 
Furthermore, our proposal recovers FedSGD when $\beta = 0$ and FedAvg when $\beta = 1$, see Equation~\ref{eq: theta-decay}. 
Viewing our proposal as a function of $t = nK + k$ iterations, it can be thought of as the specific locally decaying, cyclic learning rate where $\eta_t = \eta \times \beta^{(t \bmod K)}$.
\begin{align}
    & \theta_i^{(n, k)} = \theta_i^{(n, k - 1)} - \eta \beta^{(k - 1)} \times g_i^{(n, k - 1)} \textrm{ for } k = 1, \ldots, K \textrm{ where } \beta \in \left[0, 1\right] \\
    & g_{Proposal} = \dfrac{\theta_i^{(n, K)} - \theta_i^{(n, 0)}}{-\eta} = \sum_{j = 1}^K \left( \dfrac{\theta_i^{(n, j)} - \theta_i^{(n, j - 1)}}{\eta} \right) = \sum_{j = 1}^K \beta^{(j - 1)} g_i^{(n, j)} \label{eq: theta-decay} 
\end{align}


\section{Summary}

We suggest a locally decaying, cyclic learning rate for better balancing initial model success and personalization.
We plan to extend the work of \cite{fedavg-conv} to our learning rate schedule to how our learning rate scheduling impacts the convergence rate under non-independent and non-identically distributed local data sets.
Next, we will produce experimental results across vision, text, and graph benchmark data sets, to evaluate the effectiveness of our proposal regarding various generalization and fairness metrics.
In summary, we want to answer whether local learning rate scheduling can improve the generalization and fairness of Personalized Federated Learning without additional computation or memory cost. 


%
{  
    \small 
    \bibliographystyle{iclr2023_conference}
    \bibliography{references}
}


%
%\clearpage
%\appendix

\end{document}